\documentclass[10pt]{article}

\usepackage[margin=1in]{geometry}
\usepackage{amsmath}
\usepackage{booktabs}
\usepackage{graphicx}
\usepackage{hyperref}
\usepackage{amssymb}
\usepackage{multirow}
\usepackage[section]{placeins}

\title{Protocol-Sensitive Component Evaluation of BNNeck and ArcFace for Cross-Session Palmprint Recognition}
\author{PALM\_CGK\_BASE Project Team}
\date{\today}

\begin{document}

\maketitle

\begin{abstract}
Evaluation protocols can reverse conclusions about metric-learning components in contactless palmprint recognition.
Seen-identity diagnostics are therefore insufficient evidence for cross-session held-out palm-class deployment.
Using Tongji as the primary benchmark, we evaluate bidirectional S1$\rightarrow$S2 and S2$\rightarrow$S1 cross-session matching over three seeds with a controlled ResNet18 matrix spanning cross-entropy, supervised contrastive learning, BNNeck, ArcFace, and their combinations under a palm-class-disjoint protocol.
Under this stricter protocol, B6 (BNNeck+ArcFace) does not improve B1 (cross-entropy+SupCon) overall, and the results show directional asymmetry rather than uniform component gains.
The paper contributes an audited, reproducible protocol-sensitive evaluation rather than a new architecture or state-of-the-art claim; B5 (BNNeck+cross-entropy) is reported only as the highest observed mean on selected strict Tongji metrics, not as statistically established superiority.
\end{abstract}



\section{Introduction}

Contactless palmprint recognition is a biometric matching problem in which both identification accuracy and low-false-accept verification behavior matter. In practice, a system may be required to retrieve the correct gallery identity while also operating at strict false-acceptance rates under changes in illumination, pose, deformation, region-of-interest quality, and acquisition session. These requirements make the evaluation protocol a central part of the scientific claim. A component that improves performance when development and test identities overlap may not preserve the same ranking when the test palms are held out, the gallery and probe come from different sessions, and the operating point is constrained by a low empirical false-accept rate.

Recent palmprint systems increasingly reuse representation-learning components developed for face recognition and person re-identification, including angular-margin classifiers, supervised contrastive losses, and normalization necks. These components are plausible choices for palmprint embeddings, but their effect is not protocol-independent. In a biometric setting, the relevant question is not only whether a component improves a diagnostic split, but whether the improvement remains under held-out identities or palms, cross-session gallery/probe construction, and conservative verification thresholds. This paper addresses that question through an audited component-level evaluation rather than by proposing a new architecture.

We ask a specific protocol-level question: do conclusions about common recognition heads and losses remain stable when evaluation moves from a seen-identity diagnostic setting to a held-out palm-class cross-session protocol? The question is intentionally narrower than state-of-the-art palmprint recognition. Its importance is methodological: if component rankings change under different split and threshold policies, then claims about recognition-head improvements must be tied to the protocol under which they were measured.

This paper makes four contributions:
\begin{enumerate}
    \item It defines and audits a held-out palm-class Tongji cross-session protocol with bidirectional S1$\rightarrow$S2 and S2$\rightarrow$S1 evaluation, explicit gallery/probe construction, split hashes, and overlap checks.
    \item It evaluates a controlled ResNet18 component matrix covering cross-entropy, supervised contrastive learning, BNNeck, angular-margin supervision (ArcFace, CosFace), and selected combinations under identical checkpoint-selection and matching policies.
    \item It reports conservative low-FAR verification metrics, including TAR@FAR computed only at thresholds whose empirical false-accept rate does not exceed the target.
    \item It analyzes protocol sensitivity: the BNNeck+ArcFace configuration that appears favorable in a seen-identity diagnostic setting does not consistently improve under held-out palm-class cross-session evaluation, and the observed component behavior differs across session direction and operating point.
\end{enumerate}
The paper is therefore a reproducible empirical evaluation, not a new-architecture or state-of-the-art performance claim.




\section{Related Work}

\subsection{Palmprint Recognition}

Palmprint recognition has historically relied on local texture structure, including line, orientation, and filter-bank representations. Recent survey evidence describes the field's transition toward deep feature learning, with learned representations increasingly used for ROI processing, feature extraction, matching, and robustness analysis \cite{gao2025deeplearning}. This paper follows the embedding-learning direction rather than presenting a new ROI detector or a handcrafted texture descriptor.

Gabor and other texture filters remain important context for palmprint recognition because they encode local ridge and crease patterns. However, the current project evidence does not support a Gabor-centered paper claim: the fixed-Gabor B2 variant is a useful baseline, but B6 provides the stronger final result on the Tongji multi-seed cross-session protocol. Accordingly, this work does not claim that Gabor filtering is the source of the final improvement.

\subsection{Metric Learning for Biometric Embeddings}

Biometric recognition systems depend on embedding geometry: genuine pairs should remain close while impostor pairs should be separated, especially at strict false-acceptance operating points. Supervised contrastive learning explicitly pulls same-class examples together and pushes different-class examples apart \cite{khosla2020supervised}. The B1 baseline in this project uses this idea together with cross entropy, forming a strong baseline for the Tongji cross-session protocol.

Margin-based normalized classifiers provide another route to discriminative embeddings. ArcFace adds an angular margin to normalized features and class weights, encouraging separation in angular space \cite{deng2019arcface}. BNNeck was introduced in person re-identification as a normalization neck that separates feature spaces used for metric comparison and classification supervision \cite{luo2019strong}. B6 adapts these two ideas to the palmprint setting: ArcFace shapes the angular classifier geometry, while BNNeck provides the post-BN embedding used for cosine verification.

\subsection{Cross-Session and Cross-Domain Robustness}

Cross-session and cross-domain robustness are central issues for deployable biometric systems. The report reviewed for this draft mentions several palmprint domain-adaptation, style-transfer, and synthetic-data works, but many of those entries still need full author-title-venue-year verification before they are added to the bibliography. For this reason, the present draft uses them only as motivation for future literature expansion, not as cited evidence.

This paper focuses on a narrower and directly measured setting: Tongji cross-session recognition with S1->S2 and S2->S1 protocols. The results should therefore be read as evidence for cross-session Tongji verification, not as a broad cross-domain benchmark claim. External datasets and verified cross-domain comparisons remain future work.




\section{Evaluated Component Matrix}

We do not introduce a new architecture. Instead, we define a controlled component matrix that isolates the effect of SupCon, ArcFace, BNNeck, and their combination under identical backbone, optimizer, embedding dimensionality, checkpoint-selection, and evaluation protocols.

\subsection{Component Matrix and Evaluation Scope}

This work evaluates a controlled family of ResNet18 embedding models rather than proposing a new palmprint architecture. All tested variants use the same backbone family and a 256-D embedding size; they differ only in the use of supervised contrastive learning, BNNeck, ArcFace, and a light SupCon term. The main variants are:
\begin{itemize}
    \item B0: ResNet18 + cross entropy;
    \item B1: ResNet18 + cross entropy + supervised contrastive learning;
    \item B4: ResNet18 + ArcFace;
    \item B5: ResNet18 + BNNeck + cross entropy;
    \item B6: ResNet18 + BNNeck + ArcFace;
    \item B7: ResNet18 + BNNeck + ArcFace + light supervised contrastive learning.
\end{itemize}
B1 is the current CE + SupCon baseline. B6 is the initially hypothesized BNNeck + ArcFace variant because it appeared favorable in the seen-identity diagnostic setting. B5 is included to isolate the effect of BNNeck without ArcFace. The strict Tongji component matrix is used to distinguish the original B6 hypothesis from the broader component behavior observed under palm-class-disjoint evaluation.

\subsection{BNNeck Embedding}

BNNeck was proposed as a normalization neck for separating the feature space used for metric comparison from the feature space optimized by classification supervision \cite{luo2019strong}. Let $h_i$ denote the output feature of the ResNet18 backbone for sample $i$. A linear projection maps this feature to a 256-D vector:
\begin{equation}
    z_i = W_e h_i + b_e.
\end{equation}
BNNeck applies batch normalization to this projected vector:
\begin{equation}
    \tilde{z}_i = \mathrm{BN}(z_i).
\end{equation}
For BNNeck variants, the evaluation embedding is the post-BN representation:
\begin{equation}
    f_i = \frac{\tilde{z}_i}{\lVert \tilde{z}_i \rVert_2}.
\end{equation}
For non-BNNeck variants, the corresponding 256-D embedding is L2-normalized before matching. Gallery and probe embeddings are compared with cosine similarity.

\subsection{ArcFace Supervision}

ArcFace uses normalized features and normalized class weights to impose an additive angular margin \cite{deng2019arcface}. Let $W_j$ be the normalized classifier weight for class $j$, and let $\theta_j$ be the angle between the normalized embedding $f_i$ and $W_j$. For the ground-truth class $y_i$, ArcFace adds an angular margin $m$ and scales logits by $s$:
\begin{equation}
    \ell_i =
    -\log
    \frac{
        \exp \left(s \cos(\theta_{y_i} + m)\right)
    }{
        \exp \left(s \cos(\theta_{y_i} + m)\right)
        + \sum_{j \ne y_i} \exp \left(s \cos(\theta_j)\right)
    }.
\end{equation}
In the current B6 configuration, $s=30.0$ and $m=0.5$. ArcFace directly shapes angular classifier geometry, but the experiments test whether that training geometry actually improves held-out palm-class verification and identification metrics.

\subsection{What Is Not Claimed as the Main Method}

Earlier Gabor variants remain useful baselines and ablations because Gabor-style orientation coding is historically important in palmprint recognition \cite{kong2004competitive,genovese2019palmnet,liang2021compnet}, but they are not the main supported claim in this paper. The present evidence does not support a fixed-Gabor, learnable-Gabor, universal BNNeck + ArcFace, state-of-the-art, or cross-dataset robustness claim. The supported claim is a scoped component evaluation under the stated Tongji and IITD protocols.




\section{Experimental Setup}

\subsection{Datasets and Protocols}

The primary evaluation uses the Tongji cross-session protocol in two directions:
\begin{itemize}
    \item S1->S2: training/gallery data from session 1 and probe data from session 2.
    \item S2->S1: training/gallery data from session 2 and probe data from session 1.
\end{itemize}
The main reported Tongji numbers are bidirectional averages over S1->S2 and S2->S1 with seeds 42, 2026, and 2705.

IITD within split is used as a secondary validation. The IITD result is treated as a scope and saturation check, not as the primary claim evidence.

\subsection{Compared Methods}

The paper compares the following methods:
\begin{itemize}
    \item B1: ResNet18 + cross entropy + supervised contrastive learning.
    \item B2: FixedGaborResNet18.
    \item B4: ResNet18 + ArcFace.
    \item B6: ResNet18 + BNNeck + ArcFace.
    \item B7: ResNet18 + BNNeck + ArcFace + light SupCon.
\end{itemize}

\subsection{Training Recipe}

The common training recipe used by the final compared runs is:
\begin{itemize}
    \item epochs: 60;
    \item learning rate: $1 \times 10^{-4}$;
    \item weight decay: $1 \times 10^{-4}$;
    \item embedding dimension: 256;
    \item Tongji seeds: 42, 2026, 2705.
\end{itemize}
B6 uses ArcFace with scale $s=30.0$ and margin $m=0.5$. B7 adds a light supervised contrastive term for ablation.

\subsection{Metrics}

The reported metrics are Rank-1, Rank-5, Macro-F1, EER, TAR@FAR=1e-2, TAR@FAR=1e-3, parameter count, FLOPs, and average inference latency. TAR@FAR=1e-3 is emphasized because it measures strict low-FAR verification behavior.

\subsection{Hardware}

The exact hardware and software environment should be filled in after a verified environment record is available. The current evidence files report latency values, but the current draft does not yet include a verified full hardware description.




\section{Results}

\subsection{Tongji Multi-Seed Results}

Table~\ref{tab:tongji_multiseed} reports bidirectional Tongji averages over seeds 42, 2026, and 2705. B6 outperforms both the B1 CE + SupCon baseline and the B2 fixed-Gabor variant on the main identification and verification metrics.

\begin{table*}[t]
\centering
\caption{Tongji bidirectional multi-seed results. Values are mean +/- standard deviation over seeds 42, 2026, and 2705.}
\label{tab:tongji_multiseed}
\resizebox{\textwidth}{!}{%
\begin{tabular}{llccccccccc}
\toprule
Method & Protocol & Rank-1 & Rank-5 & Macro-F1 & EER & TAR@1e-2 & TAR@1e-3 & Time & Params & FLOPs \\
\midrule
B1 ResNet18 + CE + SupCon & Bidir. Avg. & 94.37 +/- 0.88 & 96.24 +/- 0.67 & 93.61 +/- 1.05 & 1.95 +/- 0.30 & 96.92 +/- 0.66 & 91.56 +/- 1.91 & 1.75 +/- 0.07 ms & 11.46M & 1.82G \\
B2 FixedGaborResNet18 & Bidir. Avg. & 93.71 +/- 0.55 & 95.60 +/- 0.26 & 92.81 +/- 0.60 & 2.09 +/- 0.21 & 96.90 +/- 0.37 & 91.78 +/- 0.74 & 4.87 +/- 1.71 ms & 11.82M & 2.71G \\
B6 ResNet18 + BNNeck + ArcFace & Bidir. Avg. & 96.07 +/- 0.73 & 97.45 +/- 0.59 & 95.56 +/- 0.84 & 1.45 +/- 0.30 & 98.23 +/- 0.53 & 95.20 +/- 1.21 & 4.89 +/- 0.81 ms & 11.46M & 1.82G \\
\bottomrule
\end{tabular}%
}
\end{table*}

Relative to B1, B6 improves Rank-1 by +1.70 percentage points, reduces EER by -0.51 percentage points, and improves TAR@FAR=1e-3 by +3.63 percentage points. These gains support the scoped cross-session claim for BNNeck + ArcFace.

Relative to B2, B6 improves Rank-1 by +2.36 percentage points and TAR@FAR=1e-3 by +3.41 percentage points. B6 also uses 0.36M fewer parameters and 0.89G fewer FLOPs than B2 in the recorded summary. This makes B6 the stronger final method direction than the earlier fixed-Gabor variant.

\subsection{Direct Comparisons}

\begin{table}[t]
\centering
\caption{Bidirectional 3-seed average deltas. Positive values are better for Rank-1, Rank-5, Macro-F1, and TAR. Negative values are better for EER, time, parameters, and FLOPs.}
\label{tab:direct_comparison}
\resizebox{\columnwidth}{!}{%
\begin{tabular}{lccccccccc}
\toprule
Comparison & Rank-1 & Rank-5 & Macro-F1 & EER & TAR@1e-2 & TAR@1e-3 & Time & Params & FLOPs \\
\midrule
B6 - B1 & +1.70 & +1.21 & +1.95 & -0.51 & +1.31 & +3.63 & +3.14 ms & +0.00M & +0.00G \\
B6 - B2 & +2.36 & +1.85 & +2.75 & -0.65 & +1.33 & +3.41 & +0.03 ms & -0.36M & -0.89G \\
\bottomrule
\end{tabular}%
}
\end{table}

\subsection{IITD Secondary Validation}

Table~\ref{tab:iitd} reports the IITD within split. B6 remains competitive and improves over B2, but B1 remains the strongest method on this split. This result is important for scope control: the paper should not claim universal superiority.

\begin{table}[t]
\centering
\caption{IITD within split seed42 results.}
\label{tab:iitd}
\resizebox{\columnwidth}{!}{%
\begin{tabular}{lcccccc}
\toprule
Method & Rank-1 & Rank-5 & Macro-F1 & EER & TAR@1e-2 & TAR@1e-3 \\
\midrule
B1 ResNet18 + CE + SupCon & 99.13 & 99.57 & 98.84 & 0.36 & 99.70 & 99.41 \\
B2 FixedGaborResNet18 & 98.26 & 99.13 & 97.68 & 0.71 & 99.29 & 98.93 \\
B6 ResNet18 + BNNeck + ArcFace & 98.48 & 99.13 & 97.97 & 0.66 & 99.41 & 99.11 \\
\bottomrule
\end{tabular}%
}
\end{table}




\section{Ablation Scope}
\label{sec:ablation}

The B1, B4, B6, and B7 ablation results are used as diagnostic evidence for interpreting the effect of ArcFace, BNNeck, and light SupCon supervision. These results are not used as the final basis for a general improvement claim because the final Tongji results are reported under the stricter development/test subject-disjoint protocol.

\begin{table*}[t]
\centering
\scriptsize
\setlength{\tabcolsep}{3pt}
\caption{Diagnostic Tongji seed42 ablation under the original seen-identity protocol. These results are not the final subject-disjoint evidence.}
\label{tab:ablation_scope}
\resizebox{\textwidth}{!}{
\begin{tabular}{llcccccc}
\toprule
Method & Description & Rank-1 & Rank-5 & Macro-F1 & EER & TAR@1e-2 & TAR@1e-3 \\
\midrule
B1 & ResNet18 + CE + SupCon & 93.39 & 95.47 & 92.43 & 2.30 & 96.23 & 89.75 \\
B4 & ResNet18 + ArcFace & 94.35 & 96.40 & 93.65 & 2.03 & 96.83 & 91.86 \\
B6 & ResNet18 + BNNeck + ArcFace & 96.80 & 97.88 & 96.40 & 1.14 & 98.77 & 96.53 \\
B7 & ResNet18 + BNNeck + ArcFace + light SupCon & 96.97 & 97.98 & 96.60 & 1.13 & 98.77 & 96.50 \\
\bottomrule
\end{tabular}%
}
\end{table*}

The diagnostic ablation shows why B6 was worth evaluating under a stricter protocol. ArcFace and BNNeck improved the reported seen-identity ablation metrics, and light SupCon in B7 did not materially change the result relative to B6. However, this ablation should be read as hyperparameter and architectural evidence, not as proof that B6 is consistently superior.

The subject--disjoint Tongji evaluation in Section~\ref{sec:results} is the decisive evidence for the final paper claim. Under that protocol, B6 does not improve over B1 overall. Therefore, the ablation supports a protocol-sensitivity interpretation rather than a method-dominance claim.





\section{Discussion}

The subject-disjoint results change the interpretation of the BNNeck + ArcFace variant. In the original seen-identity Tongji protocol, B6 appears favorable. However, when development and test identities are made subject-disjoint, B6 does not outperform the B1 CE + SupCon baseline overall. This difference shows that the evaluation protocol is a first-order determinant of the empirical conclusion.

The main failure mode is not uniform across session directions. B6 has a limited gain at the strict TAR@FAR=$10^{-3}$ operating point in the S1$\rightarrow$S2 direction, but it trails B1 on most other S1$\rightarrow$S2 metrics and underperforms B1 across all reported metrics in S2$\rightarrow$S1. This indicates that angular-margin supervision and BNNeck can shape useful embedding geometry in some settings, but that the effect does not translate reliably to held-out subjects under both session directions.

The IITD within-dataset subject-disjoint protocol is prepared as a secondary validation setting, but it is not used in the current evidence package. When metrics become available, it should be read as within-dataset secondary validation, not as cross-session evidence.

These results do not imply that BNNeck or ArcFace are ineffective for palmprint embeddings. They indicate that their benefit is conditioned on protocol design, identity overlap, and session direction. A fair report must therefore distinguish seen-identity evidence from held-out subject evidence and should avoid claims of dataset-independent dominance.

External benchmark leadership and cross-dataset robustness are not claimed. Future work should evaluate additional datasets, add verified cross-dataset protocols, and test whether the direction-specific failure observed on Tongji is due to session image characteristics, split design, or training objective mismatch.





\section{Conclusion}

This study shows that conclusions about common recognition heads and losses for palmprint embeddings depend strongly on the evaluation protocol. Under a seen-identity diagnostic setting, BNNeck+ArcFace (M6) can appear highly favorable, but under held-out palm-class Tongji cross-session evaluation it does not provide a consistent improvement over a cross-entropy plus supervised-contrastive baseline (M1). The strict protocol also reveals session-direction asymmetry and differences between Rank-1 identification and low-FAR verification behavior. 

The strict component ablation further shows that M4 (BNNeck + CE) has the highest observed mean on the selected Rank-1 and low-FAR TAR criteria among the tested learned variants, indicating that while BNNeck can shape useful embedding geometry under some session directions, ArcFace-based margins do not provide a reliable additional gain under the strict palm-class-disjoint protocol.

The corrected IITD palm-class-disjoint within-session secondary validation supports this conclusion, showing a near-tie between M1 and M6 rather than a clear M6 advantage. The fixed Gabor texture reference baseline further contextualizes these results, showing that classical texture descriptors can remain competitive for nearest-neighbor Rank-1 identification while performing much worse than learned embeddings in the strict low-FAR verification regime.

These results do not imply that BNNeck, ArcFace, or supervised contrastive learning are ineffective for palmprints; rather, they show that their apparent value is conditional on split design, session direction, and threshold policy. The main contribution is an audited, reproducible evaluation framework and a calibrated empirical finding: component-level palmprint claims should be tied to held-out class protocols and conservative biometric metrics.



\bibliographystyle{IEEEtran}
\bibliography{references}

\end{document}
